% !PRINTED 439
\weberpage{439}
\begin{center}
\S.\,138.

\gesperrt{Determination of signs. Quadratic residues.}
\end{center}

In the preceding section we derived two formulae in which signs still remained to be determined, namely
\[
2^{\frac{n-1}{2}} \overset{\nu}{\Pi} \cos\frac{2\nu m\pi}{n} = \pm 1, \tag{1}
\]
\[
2^{\frac{n-1}{2}} \overset{\nu}{\Pi} \sin\frac{2\nu m\pi}{n} = \pm \sqrt{n}, \tag{2}
\]
where $m$ is relatively prime to the odd number $n$ and $\nu$ runs through the following sequence.
\[
1,\ 2,\ \cdots\ \frac{n-1}{2} \tag{$\nu$}
\]
We shall here retain the symbol $\nu$ for the numbers of this sequence.

Every integer whatsoever gives, upon division by $n$, as residue one of the numbers $\nu$ or $0$ or $n-\nu$. Instead of $n-\nu$ we may, if we also admit negative residues, choose $-\nu$; and if we accordingly set aside the numbers divisible by $n$, there remains one of the numbers $\pm\nu$ as residue. We call this the \gesperrt{absolutely least residue} (in distinction from the least residue in the ordinary sense, which is taken from the numbers $1, 2, \ldots n-1$).

Now let $m$ be a number coprime to $n$; we consider the sequence of numbers
\[
m,\ 2m,\ 3m \cdots \frac{n-1}{2}\, m \tag{$m\nu$}
\]
and form for each the absolutely least residue $\varrho$:
\[
\pm \nu_1,\ \pm \nu_2,\ \pm \nu_3 \cdots \pm \nu_{\frac{1}{2}(n-1)}; \tag{$\varrho$}
\]
among these numbers $\varrho$ no two equal ones occur, and also no two opposite ones; for if the sum or the difference of two numbers $(\varrho)$, that is $\varrho+\varrho'$ or $\varrho-\varrho'$, were equal to zero, then for two distinct numbers $\nu, \nu'$ from $(\nu)$
\[
m(\nu \pm \nu')
\]
would have to be divisible by $n$, and hence $\nu \pm \nu'$ would also have to be divisible by $n$, which is impossible, since each of the numbers $\nu, \nu'$ is less than $n:2$,
%CONT>>
